% Part: normal-modal-logic
% Chapter: filtrations
% Section: finite

\documentclass[../../../include/open-logic-section]{subfiles}

\begin{document}

\olfileid[id]{nml}{fil}{fin}

\olsection{Filtrasi Bersifat Berhingga}

Kita telah mendefinisikan filtrasi bagi setiap himpunan~$\Gamma$ yang tertutup
terhadap sub-!!{formula}s. Definisi itu sendiri sama sekali tidak menjamin
bahwa filtrasi bersifat berhingga. Bahkan, ketika $\Gamma$ tak berhingga
(misalnya, ketika himpunan itu terdiri atas semua !!{formula}s),
filtrasinya sangat mungkin tak berhingga. Namun, jika $\Gamma$ berhingga
(misalnya, ketika himpunan itu merupakan himpunan sub-!!{formula}s dari suatu
!!{formula}~$!A$ tertentu), setiap filtrasi melalui~$\Gamma$ juga berhingga.

\begin{prop}\ollabel{prop:filt-are-finite}
  Jika $\Gamma$ berhingga, maka setiap filtrasi $\mModel{M^*}$ dari suatu
  model $\mModel{M}$ melalui $\Gamma$ juga berhingga.
\end{prop}

\begin{proof}
  Ukuran $W^*$ adalah banyaknya kelas~$[w]$ yang berbeda di bawah relasi
  ekuivalensi~$\equiv$. Sebarang dua dunia $u$, $v$ dalam kelas semacam
  itu---yakni sebarang $u$ dan $v$ sedemikian sehingga~$u \equiv v$---
  bersepakat mengenai semua !!{formula}s~$!A$ dalam~$\Gamma$: bagi setiap
  $!A \in \Gamma$, $!A$ bernilai benar pada $u$ dan $v$ sekaligus, atau
  tidak bernilai benar pada keduanya. Jadi, setiap kelas~$[w]$ bersesuaian
  dengan suatu himpunan bagian dari~$\Gamma$, yaitu himpunan semua
  $!A \in \Gamma$ sedemikian sehingga $!A$ bernilai benar pada dunia-dunia
  dalam~$[w]$. Tidak ada dua kelas berbeda $[u]$ dan $[v]$ yang bersesuaian
  dengan himpunan bagian
  yang sama dari~$\Gamma$. Sebab, jika himpunan !!{formula}s yang bernilai
  benar pada $u$ sama dengan himpunan !!{formula}s yang bernilai benar pada
  $v$, maka $u$ dan $v$ bersepakat mengenai semua formula dalam~$\Gamma$, yakni
  $u \equiv v$. Namun, dengan demikian $[u] = [v]$. Jadi, terdapat fungsi
  yang !!a{injective} dari $W^*$ ke $\Pow{\Gamma}$, sehingga
  $\card{W^*} \le \card{\Pow{\Gamma}}$. Karena itu, jika $\Gamma$ memuat $n$
  formula, kardinalitas $W^*$ tidak lebih besar daripada~$2^n$.
\end{proof}

\end{document}
